%! Parametric Functions — Unit 8, Lesson 8.1 — Group Activity
\documentclass[10pt]{article}
\usepackage{precalculus-article}
\usepackage{precalculus-boxes}
\newcommand{\TallMath}[1]{$\displaystyle #1\rule[-1.4em]{0pt}{3.2em}$}

\begin{document}

\pageheader{Unit 8, Lesson 8.1}{Group Activity}
\namepartnerperiod

\vspace{0.12in}

\begin{scenariobox}[Tracking a Particle's Path]{sky}
A particle moves in the $xy$-plane. Its position at time $t$ seconds
is given by the parametric function
\[
  f(t) = \bigl(t^2,\; t - 1\bigr), \quad -2 \le t \le 3.
\]
Here $x(t) = t^2$ gives the particle's horizontal position and
$y(t) = t - 1$ gives its vertical position (both in meters).
\end{scenariobox}

\vspace{0.10in}

\begin{notesbox}{Part 1 --- Building the Table}

\textbf{1.}\quad Complete the table of values. Leave no cell blank.

\vspace{6pt}
\begin{center}
\renewcommand{\arraystretch}{1.7}
\begin{tabular}{|c|c|c|c|}
\hline
$t$ (sec) & $x(t) = t^2$ & $y(t) = t - 1$ & Position $(x,\; y)$ \\
\hline
$-2$ & \blank{1.8cm} & \blank{1.8cm} & \blank{2.6cm} \\
\hline
$-1$ & \blank{1.8cm} & \blank{1.8cm} & \blank{2.6cm} \\
\hline
$0$  & \blank{1.8cm} & \blank{1.8cm} & \blank{2.6cm} \\
\hline
$1$  & \blank{1.8cm} & \blank{1.8cm} & \blank{2.6cm} \\
\hline
$2$  & \blank{1.8cm} & \blank{1.8cm} & \blank{2.6cm} \\
\hline
$3$  & \blank{1.8cm} & \blank{1.8cm} & \blank{2.6cm} \\
\hline
\end{tabular}
\end{center}

\vspace{8pt}
\textbf{2.}\quad What is the particle's starting position (at $t = -2$)?
What is its ending position (at $t = 3$)?
\writeline

\end{notesbox}

\vspace{0.08in}

\begin{notesbox}{Part 2 --- Analyzing the Path}

\textbf{3.}\quad Look at the rows for $t = -1$ and $t = 1$.
Both give the same value of $x(t)$. What does this mean for the
particle's path? Does the particle visit the same location twice?
Explain.
\writelines{3}

\vspace{0.04in}
\textbf{4.}\quad Solve: at what time $t$ does the particle cross
the $y$-axis (where $x(t) = 0$)? Show your algebra, then state
the particle's position at that moment.
\writelines{3}

\vspace{0.04in}
\textbf{5.}\quad Solve: at what time $t$ is the particle on the
$x$-axis (where $y(t) = 0$)? Show your algebra, then state the
particle's position at that moment.
\writelines{3}

\end{notesbox}

\vspace{0.08in}

\begin{notesbox}{Part 3 --- Extension (Tier E)}

\textbf{6.}\quad From your table, you can see that $x = 1$ occurs
at both $t = -1$ and $t = 1$.
\begin{itemize}[leftmargin=1.4em, itemsep=3pt]
  \item[(a)] Write the two distinct ordered pairs $(x, y)$ where $x = 1$.
    \writeline
  \item[(b)] Can the parametric curve be expressed as a function
    $y = g(x)$ on the entire domain $-2 \le t \le 3$? Explain why
    or why not using the vertical line test.
    \writelines{3}
\end{itemize}

\end{notesbox}

\end{document}
